\uuid{tgTq}
\titre{Changement de variables}
\auteur{hanani ; mzari}
\organisation{exo7}
\licence{CC BY-NC-SA 4.0}
\datecreate{2020-01-20}
\niveau{L1}
\module{Analyse}
\chapitre{Calcul d'intégrales}
\sousChapitre{Primitives diverses}
\difficulte{2}
\theme{Primitives, changement de variables}
\exo7qcmid{lille-2.primitives.013}

\contenu{
\begin{question}
\qclassification{127.11}
\qtags{motcle=changement de variables}
Parmi les égalités suivantes, cocher celles qui sont vraies :
\begin{answers}
\good{$\displaystyle \int 4x\sqrt{2x^2+1}\mathrm{d}x=\frac{2}{3}(2x^2+1)^{3/2}+k$, où $k\in \Rr$.}
\bad{$\displaystyle \int \sqrt{2x^2+1}\mathrm{d}x=\frac{2}{3}(2x^2+1)^{3/2}+k$, où $k\in \Rr$.}
\good{$\displaystyle \int \frac{3x\mathrm{d}x}{\sqrt{3x^2+2}}=\sqrt{3x^2+2}+k$, où $k\in \Rr$.}
\bad{$\displaystyle \int \frac{\mathrm{d}x}{2\sqrt{3x^2+2}}=\sqrt{3x^2+2}+k$, où $k\in \Rr$.}
\end{answers}
\vskip3mm
\begin{explanations}
Avec $u=2x^2+1$, on a : $\mathrm{d}u=4x\mathrm{d}x$ et
$$\int 4x\sqrt{2x^2+1}\mathrm{d}x=\int \sqrt{u}\mathrm{d}u=\frac{2}{3}u^{3/2}+k=\frac{2}{3}(2x^2+1)^{3/2}+k,\; k\in \Rr.$$
De m\^eme, avec $u=3x^2+2$, on a : $\mathrm{d}u=6x\mathrm{d}x$ et
$$\int \frac{3x\mathrm{d}x}{\sqrt{3x^2+2}}=\int \frac{\mathrm{d}u}{2\sqrt{u}}=\sqrt{u}+k=\sqrt{3x^2+2}+k,\; k\in \Rr.$$
\end{explanations}
\end{question}
}
